% source: An algebraic approach to the Riemann-Roch theorem and the arithmetic theory of function fields (Athanasios Papaioannou), http://www.fen.bilkent.edu.tr/~franz/mat/Riemann-Roch.pdf
% pdf_page: 0012
% retrieved: 2026-07-26
% transcribed_by: page-transcriber (AI vision, no OCR)

\documentclass[11pt]{article}
\usepackage{amsmath,amssymb,amsthm}

% Small-caps theorem style matching the source.
\newtheoremstyle{scplain}{\topsep}{\topsep}{\itshape}{0pt}{\scshape}{.}{ }{\thmname{#1} \thmnote{#3}}
\theoremstyle{scplain}
\newtheorem*{lemma}{Lemma}
\newtheorem*{definition}{Definition}
\newtheorem*{theorem}{Theorem}

% Small-caps "Proof." to match the source.
\makeatletter
\renewenvironment{proof}[1][\proofname]{%
  \par\pushQED{\qed}\normalfont\topsep6\p@\@plus6\p@\relax
  \trivlist\item[\hskip\labelsep\scshape #1\@addpunct{.}]\ignorespaces}%
  {\popQED\endtrivlist\@endpefalse}
\makeatother
\renewcommand{\qedsymbol}{$\square$}

\DeclareMathOperator{\Div}{Div}
\DeclareMathOperator{\Class}{Class}

\begin{document}

\noindent
$\deg(D_2+B)+1-g-(\deg B+g-1)
= \deg B+(\deg D_1+\deg D_2+3(1-g))$. The term in brackets is independent of
$B$ and therefore
$\dim U_1+\dim U_2-\dim\Omega_{\mathbb{F}}(-B)>0$ if the degree of $B$ is
sufficiently large. By our little fact from linear algebra, it follows that
$U_1 \cap U_2 \neq \varnothing$, as desired.

After we accomplished the above, the proof of our initial claim is easy to come:
pick $x_1 \in L(D_1+B)$, $x_2 \in L(D_2+B)$ such that
$x_1\omega_1=x_2\omega_2 \neq 0$; then
$\omega_2=(x_1x_2^{-1})\omega_1$, as desired.\hfill$\square$

We will now proceed to attach a natural divisor to any Weil differential of
$\mathbb{F}/\mathbb{K}$. Before that, we have to prove the following

\begin{lemma}[1.24]
Let $\omega$ be a non-trivial differential of $\mathbb{F}/\mathbb{K}$. Then,
there is a uniquely determined divisor
\[
  W \in M(\omega) =
    \{D \in \Div(\mathbb{F}) \mid
      \omega\big|_{A_{\mathbb{F}}(D)+\mathbb{F}} \equiv 0\}
\]
such that $D \leq W$ for any $D \in M(\omega)$.
\end{lemma}

\begin{proof}
By Riemann's inequality, there exists a constant $c$ dependent only on the
function field $\mathbb{F}/\mathbb{K}$ such that $i(D)=0$ for all
$D \in \Div(\mathbb{F})$ with the property $\deg D \geq c$ (see our remarks
following the introduction of Riemann's inequality). Since
$\dim(A_{\mathbb{F}}/(A_{\mathbb{F}}(D)+\mathbb{F}))=i(D)$, we have
$\deg D<c$ for all $D \in M(\omega)$. Thus, we may pick a divisor $W$ of maximal
degree among those.

Suppose that $W$ does not satisfy the desired property; then there exists a
divisor $D_0 \in M(\omega)$ such that $v_Q(D_0)>v_Q(W)$ at some place $Q$ of
$\mathbb{F}/\mathbb{K}$. We claim that $W+Q \in M(\omega)$, which will be a
contradiction to the maximality of $W$. Indeed, consider an adele
$\alpha=(\alpha)_P \in A_{\mathbb{F}}(W+Q)$. We may write
$\alpha=\alpha'+\alpha''$ with
\[
  \alpha'_P=(1-\delta_{PQ})\alpha_P
\]
and
\[
  \alpha''_P=\delta_{PQ}\alpha_P,
\]
where the $\delta$ above is merely the Kronecker delta. Then,
$\alpha' \in A_{\mathbb{F}}(W)$, $\alpha'' \in A_{\mathbb{F}}(D_0)$, hence
$\omega(\alpha)=\omega(\alpha')+\omega(\alpha'')=0$. Therefore, $\omega$
vanishes on $A_{\mathbb{F}}(W+Q)+\mathbb{F}$, and our claim is true. This
completes the proof of the lemma.
\end{proof}

In light of the above lemma, the following definition is legitimate

\begin{definition}[1.14]
The divisor $(\omega)$ of a Weil differential $\omega \neq 0$ is the unique
divisor $D$ of $\mathbb{F}/\mathbb{K}$ such that
\begin{enumerate}
\item[(a)] $\omega$ vanishes on $A_{\mathbb{F}}(D)+\mathbb{F}$.

\item[(b)] If $\omega$ vanishes on $A_{\mathbb{F}}(D')+\mathbb{F}$, then
$D' \leq D=(\omega)$.
\end{enumerate}

Now for any place $P$ and any $\omega,(\omega)$ as above, we define
$v_P(\omega)=v_P((\omega))$. Moreover, a place $P$ is called a zero of $\omega$
if $v_P(\omega)>0$ and a pole of $\omega$ if $v_P(\omega)>0$. A differential
$\omega$ is called regular at a place $P$ if $v_P(\omega) \geq 0$, and $\omega$
is called regular or holomorphic, if it is regular at all places
$P \in \mathbf{P}_{\mathbb{F}}$.
\end{definition}

Note now that
$\Omega_{\mathbb{F}}(D)=\{\omega \in \Omega_{\mathbb{F}} \mid
\omega=0 \text{ or } (\omega)\geq D\}$ and
$\Omega_{\mathbb{F}}(0)=\{\omega \in \Omega_{\mathbb{F}} \mid
\omega \text{ is regular}\}$. As a consequence of the fact that
$\Omega_{\mathbb{F}}(D)=i(D)$ and the definition, we obtain
$\dim\Omega_{\mathbb{F}}(0)=g$.

We will now need the following

\begin{theorem}[1.25]
\begin{enumerate}
\item[(a)] For a non-zero element $x \in \mathbb{F}$, and a non-zero
differential $\omega \in \Omega_{\mathbb{F}}$, we have
$(x\omega)=(x)+(\omega)$.

\item[(b)] Any two canonical divisors are equivalent.
\end{enumerate}
\end{theorem}

\medskip\noindent{\scshape Proof.} A simple consequence of this statement is
that the canonical divisors of $\mathbb{F}/\mathbb{K}$ form a class $[W]$ in the
divisor class group $\Class(\mathbb{F})$; this divisor class is called the
\emph{canonical class} of $\mathbb{F}/\mathbb{K}$. We now proceed to prove the
statement.

\begin{enumerate}
\item[(a)] If $\omega$ vanishes on $A_{\mathbb{F}}(D)+\mathbb{F}$, then
$x\omega$ vanishes on $A_{\mathbb{F}}(D+(x))+\mathbb{F}$, as mentioned before,
and therefore $(\omega)+(x) \leq (x\omega)$. Likewise,
$(x\omega)+(x^{-1}) \leq (\omega)$. Combining these inequalities, we obtain
$(\omega)+(x) \leq (x\omega) \leq -(x^{-1})+(\omega)=(\omega)+(x)$, which
obviously yields the desired equality.

\item[(b)] This second assertion follows from part (a) above and the fact that
$\dim_{\mathbb{F}}\Omega_{\mathbb{F}}=1$.
\end{enumerate}

\begin{center}12\end{center}

\end{document}
